% !TeX root = main.tex
\chapter{Apache Spark}
\label{ch:spark}

La thèse de référence consacre son chapitre~3 à décrire Flink \emph{avant}
d'en présenter l'implémentation (chapitre~4), parce que le dataflow proposé
est directement contraint par ce que le modèle d'exécution sait faire : état
par \emph{subtask}, \emph{watermarks}, diffusion (\emph{broadcast}). Nous
suivons la même logique avec Spark : ce chapitre décrit les mécanismes sur
lesquels s'appuie notre implémentation (chapitre~\ref{ch:implementation}),
avant de la présenter.

\section{RDD et DataFrame}

Spark expose deux abstractions principales pour les données distribuées.
Un \texttt{RDD} (\emph{Resilient Distributed Dataset}) est une collection
immuable, partitionnée, d'objets Scala arbitraires --- c'est l'abstraction
que nous utilisons pour manipuler des \texttt{streamkm.core.Point}, une
classe qui n'a pas de schéma de colonnes et ne bénéficie donc d'aucune
optimisation particulière à être représentée en \texttt{DataFrame}. Un
\texttt{DataFrame} est une table distribuée typée par un schéma de colonnes,
optimisée par le moteur Catalyst (plan logique, plan physique, génération de
code) --- c'est l'abstraction que nous utilisons pour la source du flux
(\texttt{streamkm.spark.StreamKMPlusPlus.run}), avant de repasser en RDD dès
que le traitement touche à la logique de \textsc{StreamKM++} elle-même.

Ce choix --- \texttt{DataFrame} en entrée, \texttt{RDD} dès qu'on entre dans
notre propre logique --- est délibéré : \texttt{Point} n'est pas une
\texttt{case class} (cf. \texttt{streamkm.core.Point}), et lui construire un
\texttt{Encoder} Spark pour bénéficier de Catalyst sur ce type n'aurait
apporté aucun gain (aucune requête, aucune agrégation column-oriented n'est
faite sur un \texttt{Point}), au prix d'un couplage supplémentaire entre
\texttt{core} et Spark --- exactement ce que la « règle d'or » de notre
architecture (§\ref{ch:introduction}) interdit.

\section{Micro-batch et vrai streaming}

Flink est un moteur de \emph{vrai streaming} : chaque événement est traité
dès son arrivée, avec un état persistant entre événements. Spark Structured
Streaming, malgré son nom, reste fondamentalement un moteur de
\emph{micro-batch} : le flux est découpé en petits lots (déclenchés par un
\texttt{Trigger}, par exemple un intervalle de temps ou l'arrivée d'un
fichier), et chaque lot est traité comme un \texttt{DataFrame} borné
ordinaire, via \texttt{foreachBatch}. Cette différence de modèle a une
conséquence directe et positive pour notre portage : la thèse doit détecter
elle-même la fin d'un fichier en interceptant le mécanisme de
\emph{watermark} de Flink (opérateur \texttt{HandleWatermark}, §4.1 de la
thèse) ; avec Spark, la frontière entre deux lots est \textbf{explicite} ---
\texttt{foreachBatch} est appelé une fois par micro-batch, sans bricolage
nécessaire. C'est le premier exemple concret, développé au
chapitre~\ref{ch:implementation}, de ce qui se simplifie dans le portage.

En contrepartie, le micro-batch introduit une latence structurelle
qu'un vrai moteur de streaming n'a pas : un point n'est traité qu'à la
clôture de son lot, jamais immédiatement à son arrivée. Nous revenons sur
ce compromis au chapitre~\ref{ch:discussion}.

\section{Diffusion (\texttt{broadcast})}

Quand une valeur doit être lue, identique, par toutes les tâches d'un job
(par exemple un jeu de centres déjà calculé), l'envoyer une fois par tâche
gaspillerait de la bande passante réseau proportionnellement au nombre de
partitions. \texttt{SparkContext.broadcast} envoie la valeur une seule fois
par exécuteur (et non par tâche), et la met en cache localement pour toute
la durée de vie du \texttt{Broadcast}. Nous l'utilisons dans
\texttt{streamkm.spark.CostEvaluator.cost} pour diffuser les centres avant
de calculer, en parallèle, le coût de chaque partition
(chapitre~\ref{ch:implementation}, §\ref{sec:impl-evaluation}) --- l'exact
équivalent Spark de l'opérateur \texttt{FlatMapToCentroid} + \texttt{.broadcast()}
de la thèse (§4.2).

\section{\texttt{treeReduce} et \texttt{treeAggregate}}
\label{sec:spark-treereduce}

Réduire les résultats de $N$ partitions vers le \emph{driver} peut se faire
de deux façons : \textbf{linéairement} (chaque partition envoie son résultat
directement au driver, qui les combine un par un --- $O(N)$ étapes
séquentielles sur le driver), ou en \textbf{arbre}
(\texttt{treeReduce}/\texttt{treeAggregate} : les résultats sont combinés
par paires à travers plusieurs paliers intermédiaires, $O(\log N)$ étapes
séquentielles, le reste en parallèle). Pour un grand nombre de partitions, la
réduction linéaire concentre tout le travail de fusion sur un seul n\oe{}ud ---
exactement le défaut que la thèse identifie dans son propre opérateur
\texttt{FlatMapToFinalCoreset} (§4.1, non-parallèle par construction) comme
le goulot qui plafonne le speedup au-delà de 8 c\oe{}urs (§5.3.1). C'est le point
central de notre apport par rapport à la thèse
(chapitre~\ref{ch:implementation}, §\ref{sec:impl-dataflow}) : nous
remplaçons cette fusion séquentielle par une fusion en arbre.

Nous utilisons \texttt{treeReduce} plutôt que \texttt{treeAggregate} pour la
fusion des coresets : \texttt{treeAggregate(zeroValue)(seqOp, combOp, depth)}
clone un \texttt{zeroValue} \emph{unique} (donc une seule graine aléatoire)
pour chaque partition, ce qui casserait la reproductibilité par partition
requise par notre protocole expérimental (chapitre~\ref{ch:evaluation},
§\ref{sec:eval-scalabilite}). \texttt{mapPartitionsWithIndex} (une graine
dérivée de l'indice de partition) suivi de \texttt{treeReduce} réalise la
même opération --- un accumulateur construit une fois par partition, fusionné
en arbre --- sans ce biais. Ce choix, documenté dans le code
(\texttt{StreamKMPlusPlus.mergeCoresets}), est un écart assumé par rapport à
l'exemple de pipeline de notre spécification interne, pas par rapport à ses
invariants.

\section{Partitions et c\oe{}urs : deux leviers découplés}
\label{sec:spark-partitions-coeurs}

Chez Flink, le \emph{parallélisme} est un seul nombre : il fixe à la fois le
nombre de \emph{subtasks} et le nombre de listes de \emph{buckets}
merge-and-reduce actives. Chez Spark, ces deux choses sont
\textbf{découplées} : le \textbf{nombre de partitions} détermine le nombre
d'accumulateurs \texttt{BucketManager} indépendants (donc la
\emph{qualité} --- plus de partitions, plus de coresets partiels, un
raffinement plus fin) ; le \textbf{nombre de c\oe{}urs} détermine le
\emph{débit} de traitement (donc le \emph{temps}). La thèse observe que le
coût \emph{diminue} quand le parallélisme Flink augmente (§5.3.1,
fig.~5.9) et l'attribue à la multiplication des listes de buckets --- mais
elle ne peut pas isoler cet effet de celui du nombre de c\oe{}urs, puisque les
deux varient ensemble chez Flink. Le découplage de Spark permet de tester
cette explication directement, en fixant les c\oe{}urs et en faisant varier
uniquement les partitions (chapitre~\ref{ch:evaluation},
§\ref{sec:eval-scalabilite}, expérience E3bis) --- une vérification que le
modèle de Flink ne permet tout simplement pas de faire, et qui constitue
l'une de nos contributions propres à ce rapport. En pratique, cela signifie
que \texttt{spark.sql.shuffle.partitions} et le nombre de partitions des RDD
manipulés sont toujours fixés explicitement dans notre code
(\texttt{repartition(N)}), jamais laissés à la décision de Spark.
